\subsection{Vectorization application}
The application of vectorization was carried out by replacing the old \texttt{ant.cpp} file, which contained a vector storing each \texttt{Ant} class (with its respective properties), with several separate vectors representing each of these properties.
That is, the \texttt{seed}, \texttt{state}, and position (\texttt{x} and \texttt{y}) properties of each ant were transferred to four vectors within a larger class called \texttt{Population}, whose function is to manage the change in the ants' positions.

For this adaptation, the \texttt{advance} function was modified to receive (in addition to the pheromone and terrain conditions, etc.) the index of each ant in the \texttt{Population} class vectors. It maintains the previous logic but uses the name \texttt{advance\_one} for a single ant and \texttt{advance\_all} for all ants.
Finally, for organizational purposes, other functions such as \texttt{add\_ant} and \texttt{reserve} were created to add ants to the population and to initialize their vectors, respectively.
Regarding changes in code execution, the ant-by-ant iterations (\texttt{ant.advance}) are now performed using \texttt{Population.advance\_all} with the respective indices for every ant. Table~\ref{tab:vectorization_summary} summarizes the main vectorized components.

\begin{table}[htbp]
\centering
\caption{Summary of vectorized components in the \texttt{Population} structure.}
\label{tab:vectorization_summary}
\small
\begin{tabular}{p{0.32\linewidth} p{0.60\linewidth}}
\hline
\textbf{Component} & \textbf{Brief description} \\
\hline
\texttt{Population} & Class storing ant properties in separate vectors. \\
\shortstack[l]{\texttt{Population::}\\\texttt{advance\_one}} & Advances a single ant using its vector index. \\
\shortstack[l]{\texttt{Population::}\\\texttt{advance\_all}} & Iterates over all indices to advance the population. \\
\shortstack[l]{\texttt{Population::}\\\texttt{add\_ant}} & Adds an ant to the population. \\
\shortstack[l]{\texttt{Population::}\\\texttt{reserve}} & Initializes and allocates memory for the vectors. \\
\hline
\end{tabular}
\end{table}
